% !TEX program = xelatex
% =====================================================================
%  PhotoLib 校区负责人使用指南
%  编译：xelatex 校区负责人使用指南.tex（连续两次以生成目录）
% =====================================================================
\documentclass[cn,11pt,green]{elegantbook}

% =====================================================================
%  elegant-style.tex —— PhotoLib 使用指南三份文档共用的样式与宏定义
%  由 管理员/部长/校区负责人 三份指南通过 % =====================================================================
%  elegant-style.tex —— PhotoLib 使用指南三份文档共用的样式与宏定义
%  由 管理员/部长/校区负责人 三份指南通过 % =====================================================================
%  elegant-style.tex —— PhotoLib 使用指南三份文档共用的样式与宏定义
%  由 管理员/部长/校区负责人 三份指南通过 \input{elegant-style} 引入。
%  依赖：ElegantBook 文档类（须以 XeLaTeX 编译）。
% =====================================================================

% ---- 颜色（与 ElegantBook green 主题协调）----
\definecolor{plMain}{RGB}{0,120,80}      % 主色：墨绿
\definecolor{plInk}{RGB}{38,50,56}       % 正文强调
\definecolor{admTipBg}{RGB}{236,247,240}
\definecolor{admTipFr}{RGB}{34,139,88}
\definecolor{admNoteBg}{RGB}{232,240,254}
\definecolor{admNoteFr}{RGB}{25,103,210}
\definecolor{admWarnBg}{RGB}{253,237,237}
\definecolor{admWarnFr}{RGB}{198,40,40}
\definecolor{admStepBg}{RGB}{245,243,255}
\definecolor{admStepFr}{RGB}{103,58,183}

% ---- tcolorbox（ElegantBook 已载入，这里补声明所需库）----
\usepackage{tcolorbox}
\tcbuselibrary{skins,breakable}

% 统一的四种“提示框”环境，均支持可选自定义标题：
%   \begin{tipbox}[自定义标题] ... \end{tipbox}
\newtcolorbox{tipbox}[1][小贴士]{%
  enhanced, breakable,
  colback=admTipBg, colframe=admTipFr, coltitle=white,
  fonttitle=\bfseries\sffamily, title={#1},
  boxrule=0.6pt, arc=3pt, left=8pt, right=8pt, top=6pt, bottom=6pt}

\newtcolorbox{notebox}[1][说明]{%
  enhanced, breakable,
  colback=admNoteBg, colframe=admNoteFr, coltitle=white,
  fonttitle=\bfseries\sffamily, title={#1},
  boxrule=0.6pt, arc=3pt, left=8pt, right=8pt, top=6pt, bottom=6pt}

\newtcolorbox{warnbox}[1][注意]{%
  enhanced, breakable,
  colback=admWarnBg, colframe=admWarnFr, coltitle=white,
  fonttitle=\bfseries\sffamily, title={#1},
  boxrule=0.6pt, arc=3pt, left=8pt, right=8pt, top=6pt, bottom=6pt}

\newtcolorbox{stepbox}[1][操作步骤]{%
  enhanced, breakable,
  colback=admStepBg, colframe=admStepFr, coltitle=white,
  fonttitle=\bfseries\sffamily, title={#1},
  boxrule=0.6pt, arc=3pt, left=8pt, right=8pt, top=6pt, bottom=6pt}

% ---- 表格 ----
\usepackage{booktabs}
\usepackage{array}
\usepackage{longtable}
\usepackage{makecell}
\renewcommand{\arraystretch}{1.28}

% 左对齐、可换行的定宽列
\newcolumntype{L}[1]{>{\raggedright\arraybackslash}p{#1}}

% ---- 行内“键位/菜单/字段”标记 ----
\usepackage{relsize}
% 菜单路径：系统管理 \menu{数据备份}
\newcommand{\menu}[1]{{\sffamily\bfseries\color{plMain}#1}}
% 界面按钮/字段名
\newcommand{\ui}[1]{\textsf{「#1」}}
% 权限码等等宽标识；在每个下划线后允许换行，避免长标识顶出表格边界
\newcommand{\code}[1]{{\ttfamily\smaller[0.5]%
  \renewcommand{\_}{\textunderscore\penalty0\hspace{0pt}}#1}}
% 角色徽标
\newcommand{\role}[1]{{\sffamily\bfseries\color{admNoteFr}#1}}

% ---- 超链接 ----
\hypersetup{
  colorlinks=true,
  linkcolor=plMain,
  urlcolor=admNoteFr,
  citecolor=plMain}
 引入。
%  依赖：ElegantBook 文档类（须以 XeLaTeX 编译）。
% =====================================================================

% ---- 颜色（与 ElegantBook green 主题协调）----
\definecolor{plMain}{RGB}{0,120,80}      % 主色：墨绿
\definecolor{plInk}{RGB}{38,50,56}       % 正文强调
\definecolor{admTipBg}{RGB}{236,247,240}
\definecolor{admTipFr}{RGB}{34,139,88}
\definecolor{admNoteBg}{RGB}{232,240,254}
\definecolor{admNoteFr}{RGB}{25,103,210}
\definecolor{admWarnBg}{RGB}{253,237,237}
\definecolor{admWarnFr}{RGB}{198,40,40}
\definecolor{admStepBg}{RGB}{245,243,255}
\definecolor{admStepFr}{RGB}{103,58,183}

% ---- tcolorbox（ElegantBook 已载入，这里补声明所需库）----
\usepackage{tcolorbox}
\tcbuselibrary{skins,breakable}

% 统一的四种“提示框”环境，均支持可选自定义标题：
%   \begin{tipbox}[自定义标题] ... \end{tipbox}
\newtcolorbox{tipbox}[1][小贴士]{%
  enhanced, breakable,
  colback=admTipBg, colframe=admTipFr, coltitle=white,
  fonttitle=\bfseries\sffamily, title={#1},
  boxrule=0.6pt, arc=3pt, left=8pt, right=8pt, top=6pt, bottom=6pt}

\newtcolorbox{notebox}[1][说明]{%
  enhanced, breakable,
  colback=admNoteBg, colframe=admNoteFr, coltitle=white,
  fonttitle=\bfseries\sffamily, title={#1},
  boxrule=0.6pt, arc=3pt, left=8pt, right=8pt, top=6pt, bottom=6pt}

\newtcolorbox{warnbox}[1][注意]{%
  enhanced, breakable,
  colback=admWarnBg, colframe=admWarnFr, coltitle=white,
  fonttitle=\bfseries\sffamily, title={#1},
  boxrule=0.6pt, arc=3pt, left=8pt, right=8pt, top=6pt, bottom=6pt}

\newtcolorbox{stepbox}[1][操作步骤]{%
  enhanced, breakable,
  colback=admStepBg, colframe=admStepFr, coltitle=white,
  fonttitle=\bfseries\sffamily, title={#1},
  boxrule=0.6pt, arc=3pt, left=8pt, right=8pt, top=6pt, bottom=6pt}

% ---- 表格 ----
\usepackage{booktabs}
\usepackage{array}
\usepackage{longtable}
\usepackage{makecell}
\renewcommand{\arraystretch}{1.28}

% 左对齐、可换行的定宽列
\newcolumntype{L}[1]{>{\raggedright\arraybackslash}p{#1}}

% ---- 行内“键位/菜单/字段”标记 ----
\usepackage{relsize}
% 菜单路径：系统管理 \menu{数据备份}
\newcommand{\menu}[1]{{\sffamily\bfseries\color{plMain}#1}}
% 界面按钮/字段名
\newcommand{\ui}[1]{\textsf{「#1」}}
% 权限码等等宽标识；在每个下划线后允许换行，避免长标识顶出表格边界
\newcommand{\code}[1]{{\ttfamily\smaller[0.5]%
  \renewcommand{\_}{\textunderscore\penalty0\hspace{0pt}}#1}}
% 角色徽标
\newcommand{\role}[1]{{\sffamily\bfseries\color{admNoteFr}#1}}

% ---- 超链接 ----
\hypersetup{
  colorlinks=true,
  linkcolor=plMain,
  urlcolor=admNoteFr,
  citecolor=plMain}
 引入。
%  依赖：ElegantBook 文档类（须以 XeLaTeX 编译）。
% =====================================================================

% ---- 颜色（与 ElegantBook green 主题协调）----
\definecolor{plMain}{RGB}{0,120,80}      % 主色：墨绿
\definecolor{plInk}{RGB}{38,50,56}       % 正文强调
\definecolor{admTipBg}{RGB}{236,247,240}
\definecolor{admTipFr}{RGB}{34,139,88}
\definecolor{admNoteBg}{RGB}{232,240,254}
\definecolor{admNoteFr}{RGB}{25,103,210}
\definecolor{admWarnBg}{RGB}{253,237,237}
\definecolor{admWarnFr}{RGB}{198,40,40}
\definecolor{admStepBg}{RGB}{245,243,255}
\definecolor{admStepFr}{RGB}{103,58,183}

% ---- tcolorbox（ElegantBook 已载入，这里补声明所需库）----
\usepackage{tcolorbox}
\tcbuselibrary{skins,breakable}

% 统一的四种“提示框”环境，均支持可选自定义标题：
%   \begin{tipbox}[自定义标题] ... \end{tipbox}
\newtcolorbox{tipbox}[1][小贴士]{%
  enhanced, breakable,
  colback=admTipBg, colframe=admTipFr, coltitle=white,
  fonttitle=\bfseries\sffamily, title={#1},
  boxrule=0.6pt, arc=3pt, left=8pt, right=8pt, top=6pt, bottom=6pt}

\newtcolorbox{notebox}[1][说明]{%
  enhanced, breakable,
  colback=admNoteBg, colframe=admNoteFr, coltitle=white,
  fonttitle=\bfseries\sffamily, title={#1},
  boxrule=0.6pt, arc=3pt, left=8pt, right=8pt, top=6pt, bottom=6pt}

\newtcolorbox{warnbox}[1][注意]{%
  enhanced, breakable,
  colback=admWarnBg, colframe=admWarnFr, coltitle=white,
  fonttitle=\bfseries\sffamily, title={#1},
  boxrule=0.6pt, arc=3pt, left=8pt, right=8pt, top=6pt, bottom=6pt}

\newtcolorbox{stepbox}[1][操作步骤]{%
  enhanced, breakable,
  colback=admStepBg, colframe=admStepFr, coltitle=white,
  fonttitle=\bfseries\sffamily, title={#1},
  boxrule=0.6pt, arc=3pt, left=8pt, right=8pt, top=6pt, bottom=6pt}

% ---- 表格 ----
\usepackage{booktabs}
\usepackage{array}
\usepackage{longtable}
\usepackage{makecell}
\renewcommand{\arraystretch}{1.28}

% 左对齐、可换行的定宽列
\newcolumntype{L}[1]{>{\raggedright\arraybackslash}p{#1}}

% ---- 行内“键位/菜单/字段”标记 ----
\usepackage{relsize}
% 菜单路径：系统管理 \menu{数据备份}
\newcommand{\menu}[1]{{\sffamily\bfseries\color{plMain}#1}}
% 界面按钮/字段名
\newcommand{\ui}[1]{\textsf{「#1」}}
% 权限码等等宽标识；在每个下划线后允许换行，避免长标识顶出表格边界
\newcommand{\code}[1]{{\ttfamily\smaller[0.5]%
  \renewcommand{\_}{\textunderscore\penalty0\hspace{0pt}}#1}}
% 角色徽标
\newcommand{\role}[1]{{\sffamily\bfseries\color{admNoteFr}#1}}

% ---- 超链接 ----
\hypersetup{
  colorlinks=true,
  linkcolor=plMain,
  urlcolor=admNoteFr,
  citecolor=plMain}


\title{PhotoLib 校区负责人使用指南}
\subtitle{接单 · 上传 · 交付 · 报工时}
\author{PhotoLib 项目组}
\institute{校园摄影部图片工作站}
\version{1.0}
\date{2026 年 9 月}
\extrainfo{本指南面向校区负责人（CAMPUS\_MANAGER），一步一步教你\\接需求、传图片、交作业、报工时，看完就能干活。}
\cover{cover.png}

\begin{document}

\maketitle
\frontmatter
\tableofcontents
\mainmatter

% =====================================================================
\chapter{开始之前}
% =====================================================================

\begin{introduction}
  \item 你在系统里负责什么
  \item 怎么登录、界面长什么样
\end{introduction}

\section{你的角色}

你是\textbf{校区负责人}（系统里叫 \role{CAMPUS\_MANAGER}）。你负责\textbf{本校区}的具体拍摄工作，主要做四件事：

\begin{enumerate}
  \item \textbf{接需求}：接下部长发到你校区的拍摄任务。
  \item \textbf{传图片}：把拍好的照片上传进系统。
  \item \textbf{交作业}：把图片交付给对应的需求。
  \item \textbf{报工时}：填报自己参与任务的工时。
\end{enumerate}

此外你还要\textbf{维护本校区的通讯录}（成员名单）。

\begin{notebox}[你只管自己的校区]
  你看到和能操作的，基本都限于\textbf{你被授权的校区}以及\textbf{你本人参与的任务}。
  别的校区的需求、别人上传的图片，默认你是看不到、也管不了的。这是系统的设计，不是出了问题。
\end{notebox}

\section{登录与第一次改密}

\begin{stepbox}[第一次登录]
  \begin{enumerate}
    \item 打开管理员给你的网址，输入\textbf{用户名（或邮箱）+ 密码}。
    \item 第一次登录会\textbf{强制改密码}，设一个只有你知道的强密码。
    \item 改完就进入\ui{工作台}首页了。
  \end{enumerate}
\end{stepbox}

\begin{tipbox}[忘记密码]
  系统没有自助找回。忘了密码就找\textbf{管理员或部长}帮你重置，重置后下次登录再改一次即可。
\end{tipbox}

\section{界面导览}

左边是菜单，右上角是消息铃铛和头像。你常用的是这几个：

\begin{table}[htbp]
  \centering
  \caption{校区负责人常用菜单}
  \begin{tabular}{L{2.8cm} L{9.4cm}}
    \toprule
    \textbf{菜单} & \textbf{用来做什么} \\
    \midrule
    工作台 & 首页概览。 \\
    图片需求 & 接需求、上传交付。 \\
    图片库 & 上传图片、查看和下载。 \\
    收藏图片 & 你收藏过的图片。 \\
    工时记录 & 填报自己的工时。 \\
    通讯录 & 维护本校区成员名单。 \\
    消息中心 & 查看站内通知。 \\
    \bottomrule
  \end{tabular}
\end{table}

% =====================================================================
\chapter{接需求与交付}
% =====================================================================

\begin{introduction}
  \item 怎么接下一条需求
  \item 怎么把拍好的图交上去
  \item 接了不想干了怎么退出
\end{introduction}

\section{接受需求}

部长发布需求后，\textbf{你校区已启用的负责人会收到通知}。进入\menu{图片需求}，找到发到你校区的需求，点\ui{接受}。

\begin{itemize}
  \item 你只能接\textbf{本校区}已发布 / 已接受的需求。
  \item 一条需求\textbf{可以多个人一起接}。
  \item 同一条需求你\textbf{不能重复接}，再点会提示已接受。
\end{itemize}

\section{交付图片}

接了需求之后，就可以往这条需求里\textbf{上传并交付}图片了。

\begin{stepbox}[交付一条需求]
  \begin{enumerate}
    \item 进入你已接受的需求，进入它的\textbf{交付页面}。
    \item 上传照片（上传方法见下一章）。
    \item 确认无误后\textbf{提交}，需求状态变为「已提交」，等部长确认。
    \item 如果被部长\textbf{退回}，按意见补拍或重传，再提交一次。
  \end{enumerate}
\end{stepbox}

\begin{tipbox}[被退回是正常的]
  被退回不代表做错了，通常是需要补拍或换图。按部长的意见处理后重新提交即可。
\end{tipbox}

\section{退出需求}

如果你接了之后确实参与不了，可以\textbf{退出}这条需求。但有个前提：

\begin{warnbox}[传过图 / 报过工时就不能退了]
  只要你\textbf{已经为这条需求上传过图片，或填报过工时}，就\textbf{不能再退出}了。
  这是为了保护已经产生的数据。如果确实需要处理，联系部长或管理员。
\end{warnbox}

\begin{notebox}[最后一个人退出会怎样]
  如果一条需求所有接单人都退光了，需求会自动\textbf{退回「已发布」}状态，重新等人来接。
\end{notebox}

% =====================================================================
\chapter{上传图片}
% =====================================================================

\begin{introduction}
  \item 单张、批量、压缩包三种上传方式
  \item 「拍摄者」为什么必须从通讯录选
  \item 上传有哪些限制
\end{introduction}

\section{三种上传方式}

\begin{table}[htbp]
  \centering
  \caption{三种上传方式怎么选}
  \begin{tabular}{L{2.6cm} L{9.6cm}}
    \toprule
    \textbf{方式} & \textbf{适用场景} \\
    \midrule
    单张上传 & 只传一两张，想仔细填每张的信息。 \\
    批量文件 & 一次选多张照片一起传（最多 100 张）。 \\
    ZIP 压缩包 & 照片很多时，打成一个压缩包上传，系统自动解压。 \\
    \bottomrule
  \end{tabular}
\end{table}

\begin{notebox}[上传的限制]
  \begin{itemize}
    \item 只支持 \textbf{JPG / JPEG / PNG}，单张最大 100 MiB。
    \item 一次批量最多 \textbf{100 张}；ZIP 压缩包最大约 1.5 GB，解压后总量最大 10 GiB。
    \item ZIP 里只会接收图片文件，其他文件、文件夹会被跳过。
  \end{itemize}
\end{notebox}

\section{拍摄者必须从通讯录选}

上传时要填\textbf{拍摄者}，它\textbf{只能从本校区通讯录里选}，不能手打名字。

\begin{warnbox}[人不在名单里怎么办]
  如果要选的拍摄者不在下拉列表里，先去\menu{通讯录}把他\textbf{加进本校区名单}，再回来上传。
  通讯录为空时，上传页也会提示你先去维护通讯录。
\end{warnbox}

\section{拍摄时间会自动读取}

系统会尝试从照片里读出\textbf{拍摄时间}（EXIF 信息）自动填好。

\begin{itemize}
  \item 单张上传：自动填这张照片的拍摄时间。
  \item 需求多图上传：统一采用\textbf{第一张照片}的拍摄时间。
  \item 读不出来时，你可以\textbf{手动选}拍摄时间。
\end{itemize}

\section{上传之后}

上传后照片会经过系统的\textbf{处理}（压缩、生成缩略图），处理好才会出现在图库里。单张上传时页面会自动等待处理完成，稍等即可。

\begin{tipbox}[重复的图传不上去]
  单张上传时，系统会检查这张图是不是\textbf{已经传过}。如果完全一样，会提示重复，不用重复传。
\end{tipbox}

% =====================================================================
\chapter{图片库与收藏}
% =====================================================================

\begin{introduction}
  \item 你在图库里能看到哪些图
  \item 怎么下载、收藏
\end{introduction}

\section{你能看到哪些图}

图库里你默认看到的是\textbf{自己上传的图片}。这是由权限决定的。

\begin{notebox}[看到的范围可能被放宽]
  管理员\textbf{可以}把你的可见范围放宽到「本校区所有人上传的图」。到底看得到多少，取决于管理员给你的设置。
  如果你觉得应该能看到同校区其他人的图却看不到，可以找管理员确认。
\end{notebox}

\section{下载与收藏}

\begin{itemize}
  \item \textbf{下载}：单张下载，或勾选多张\textbf{打包下载}（最多 200 张）。下载用的是临时安全链接，过期重新点一下即可。
  \item \textbf{收藏}：点图片上的星标收藏，之后在\menu{收藏图片}里集中看。取消收藏后它会从收藏列表里消失。
\end{itemize}

\begin{warnbox}[采用、删除不归你管]
  \textbf{标记被引（采用）是部长的权限}，你没有。图片的删除、归档等维护操作一般也\textbf{只针对你自己上传的图}，
  且已经被采用的图片\textbf{不能直接删除}。
\end{warnbox}

% =====================================================================
\chapter{填报工时}
% =====================================================================

\begin{introduction}
  \item 工时怎么填
  \item 填完之后会怎样
\end{introduction}

\section{填报}

进入\menu{工时记录}，为你\textbf{参与的任务}填报工时。

\begin{stepbox}[填一条工时]
  \begin{enumerate}
    \item 选择\textbf{工作日期}、填报的\textbf{成员}（来自本校区通讯录）、工时数等信息。
    \item 保存提交。之后由\textbf{部长或管理员确认}。
  \end{enumerate}
\end{stepbox}

\begin{notebox}[填报人来自通讯录]
  和上传拍摄者一样，工时的成员也\textbf{来自校区通讯录}。历史工时会保存当时的姓名和学号，
  日后就算把某人从通讯录删了，\textbf{也不会改动已经填好的历史工时}。
\end{notebox}

\section{确认之后}

工时提交后要等部长\textbf{确认}。只有\textbf{已确认}的工时才会计入统计。如果被\textbf{退回}，按意见修改后重新提交。

% =====================================================================
\chapter{维护通讯录}
% =====================================================================

\begin{introduction}
  \item 通讯录为什么重要
  \item 怎么增、停用、删成员
\end{introduction}

\section{为什么重要}

通讯录是\textbf{整条数据的源头}：上传时选的拍摄者、填工时时选的成员，都从这里来。名单不全，活儿就没法正常派人记账。
所以请\textbf{及时把本校区的成员维护好}。

\section{增、停用、删}

进入\menu{通讯录}（你只能维护\textbf{本校区}的名单）：

\begin{table}[htbp]
  \centering
  \caption{三种成员操作的区别}
  \begin{tabular}{L{2.2cm} L{10cm}}
    \toprule
    \textbf{操作} & \textbf{效果} \\
    \midrule
    添加 & 新增一名成员（姓名 + 学号）。 \\
    停用 & 成员保留在名单里但不再可选，\textbf{以后可以再启用}；历史记录不受影响。 \\
    删除 & \textbf{彻底移除}这一行，之后同一个学号可以\textbf{重新添加}；历史照片和工时因为存的是快照，不受影响。 \\
    \bottomrule
  \end{tabular}
\end{table}

\begin{tipbox}[停用还是删除？]
  拿不准就用\textbf{停用}——它更保险，随时能恢复。只有确定这个学号要腾出来重新用时，才用删除。
\end{tipbox}

% =====================================================================
\chapter{消息与个人设置}
% =====================================================================

\section{消息中心}

右上角\textbf{铃铛}会提示未读消息，\menu{消息中心}里能看到所有通知（比如有新需求发到你校区、工时被退回等）。
养成\textbf{每天看一眼消息}的习惯，别漏掉派给你的任务。

\section{个人设置}

点右上角\textbf{头像}可以设置头像、退出登录，也从这里改密码。

\section{遇到问题怎么办}

\begin{tipbox}[常见小问题]
  \begin{itemize}
    \item \textbf{接不了某个需求}？确认它是发到\textbf{你校区}的，且状态是已发布 / 已接受。
    \item \textbf{选不到拍摄者 / 工时成员}？先去\menu{通讯录}把人加进本校区。
    \item \textbf{图片打不开 / 下载失败}？多半是临时链接过期，刷新或重新点开即可。
    \item \textbf{退不出需求}？你已经传过图或报过工时了，找部长或管理员处理。
    \item \textbf{看不到本该看到的东西}？可能是权限范围问题，找管理员确认。
  \end{itemize}
\end{tipbox}

\begin{notebox}[记住一句话]
  你的日常就是四步循环：\textbf{接需求 → 传图片 → 交作业 → 报工时}。把通讯录维护好，其余的跟着菜单走就行。
\end{notebox}

\end{document}
